\section{Experiments}
\label{sec:experiments}

\subsection{Setup}

We evaluate on LJSpeech~\cite{ljspeech2017} with HiFi-GAN~\cite{kong2020} and
DiffWave~\cite{kong2021}: $F = 80$ mel bins, $n_{\text{fft}} = 1024$, 22.05~kHz, hop
256. \texttt{RAWMER} embeds 32 bits with $B = 40$ band bins ($20\!:\!60$),
$L = S = 8$, $N_b = 6$, $P = 6$, $C = 16$, $\alpha = 0.435$ (HiFi-GAN) / $0.355$
(DiffWave), $\Delta_{\max} = 12$, $\tau = 0.75$; code will be released. Attacks are MP3/AAC compression, $\pm 8\%$ scaling, 16~kHz resampling,
300--8000~Hz bandpass / 3~kHz lowpass filtering, additive Gaussian noise at 20/10/5~dB
SNR, and 80~ms echo. ACC is the fraction of payload bits recovered; VR the fraction of
utterances with $\text{ACC} \geq \tau$ (24 of 32 bits). PESQ and STOI are computed
against the \emph{unmarked vocoder output}, isolating embedding from synthesis cost; at
the matched point clean STOI is 0.969 for both methods on HiFi-GAN (0.960 vs.\ 0.954 on
DiffWave), so they match on intelligibility too. We use one single-speaker read-English
corpus with ground-truth mels; end-to-end TTS mels, multi-speaker or expressive speech
and re-recording are untested gaps we flag, not extrapolate over.

\begin{table*}[t]
\setlength{\belowcaptionskip}{4pt}
\centering
\footnotesize
\caption{LJSpeech~\cite{ljspeech2017}, random 2000 (seed 2026); \emph{nX} = additive
Gaussian noise at $X$\,dB SNR; ``Ref.'' marks \emph{informed} detectors, which read
the clean reference. \textbf{(a)} Bit accuracy: mel-domain methods
(\texttt{RAWMER}, MelShield~\cite{jin2026}) carry 32 bits at matched clean
PESQ\,$\approx$\,3.51, while AudioSeal~\cite{san2024} and WavMark~\cite{chen2023}
(blind, forced-decoded) carry 16 bits on random 500---context, not matched
competitors. \textbf{(b)} Verification rate at two thresholds (a clip verifies when
$\text{ACC}\!\geq\!\tau$), with FPR the chance-level rate
$\Pr[\text{Bin}(32,\tfrac12)\!\geq\!\lceil 32\tau\rceil]$: $\tau\!=\!0.61$ buys VR at
$31\times$ the FPR, so it is a bit-margin diagnostic only. Bold: best per column,
ties included.}
\label{tab:main}
\vspace{2pt}
\renewcommand{\arraystretch}{1.05}
\begin{minipage}[t]{0.465\linewidth}
\centering
{\footnotesize (a) bit accuracy\par}\vspace{2pt}
\setlength{\tabcolsep}{4pt}
\begin{tabular}{@{}l l c c c c c@{}}
\toprule
Vocoder & Method & Ref. & none & n20 & n10 & n5 \\
\midrule
\multirow{4}{*}{HiFi-GAN}
 & MelShield  & \checkmark & 0.999 & 0.903 & 0.713 & 0.633 \\
 & \texttt{RAWMER} & \checkmark & \textbf{1.000} & \textbf{0.994} & \textbf{0.941} & \textbf{0.859} \\
 & AudioSeal & --- & \textbf{1.000} & 0.974 & 0.790 & 0.647 \\
 & WavMark   & --- & \textbf{1.000} & 0.752 & 0.511 & 0.493 \\
\midrule
\multirow{4}{*}{DiffWave}
 & MelShield  & \checkmark & \textbf{1.000} & 0.969 & 0.823 & 0.722 \\
 & \texttt{RAWMER} & \checkmark & \textbf{1.000} & \textbf{0.990} & \textbf{0.922} & \textbf{0.835} \\
 & AudioSeal & --- & \textbf{1.000} & 0.966 & 0.754 & 0.618 \\
 & WavMark   & --- & \textbf{1.000} & 0.670 & 0.512 & 0.498 \\
\bottomrule
\end{tabular}
\end{minipage}\hfill
\begin{minipage}[t]{0.495\linewidth}
\centering
{\footnotesize (b) verification rate\par}\vspace{2pt}
\setlength{\tabcolsep}{4pt}
\begin{tabular}{@{}l c c l c c c@{}}
\toprule
Vocoder & $\tau$ & FPR & Method & n20 & n10 & n5 \\
\midrule
\multirow{4}{*}{HiFi-GAN}
 & \multirow{2}{*}{0.75} & \multirow{2}{*}{0.35\%} & MelShield & 0.969 & 0.434 & 0.127 \\
 & & & \texttt{RAWMER} & \textbf{1.000} & \textbf{0.991} & \textbf{0.918} \\
\cmidrule(lr){2-7}
 & \multirow{2}{*}{0.61} & \multirow{2}{*}{10.8\%} & MelShield & 0.999 & 0.874 & 0.612 \\
 & & & \texttt{RAWMER} & \textbf{1.000} & \textbf{1.000} & \textbf{0.995} \\
\midrule
\multirow{4}{*}{DiffWave}
 & \multirow{2}{*}{0.75} & \multirow{2}{*}{0.35\%} & MelShield & 0.997 & 0.841 & 0.485 \\
 & & & \texttt{RAWMER} & \textbf{1.000} & \textbf{0.983} & \textbf{0.871} \\
\cmidrule(lr){2-7}
 & \multirow{2}{*}{0.61} & \multirow{2}{*}{10.8\%} & MelShield & \textbf{1.000} & 0.973 & 0.872 \\
 & & & \texttt{RAWMER} & \textbf{1.000} & \textbf{1.000} & \textbf{0.985} \\
\bottomrule
\end{tabular}
\end{minipage}
\end{table*}

\begin{table*}[t]
\centering
\footnotesize
\caption{\texttt{RAWMER} on HiFi-GAN~\cite{kong2020}; \emph{nX} = additive Gaussian
noise at $X$\,dB SNR. \textbf{(a)} Reference compression, ACC\,/\,VR at
$\tau\!=\!0.75$ (random 2000; uint4 band is random 500); ``band'' stores the 40-bin
watermark band, ``full'' all 80 bins, and a float32 band-only reference would occupy
87.5\,KB. \textbf{(b)} Candidate count $C$, ACC (random 500); $C\!=\!1$ is
differential encoding with a random bin-group pair, i.e.\ no selection.
\textbf{(c)} Desynchronization, ACC\,/\,VR (random 500).}
\label{tab:compress}
\vspace{2pt}
\renewcommand{\arraystretch}{1.0}
\begin{minipage}[t]{0.45\linewidth}
\centering
{\footnotesize (a) reference compression\par}\vspace{2pt}
\setlength{\tabcolsep}{3pt}
\begin{tabular}{@{}l r c c c@{}}
\toprule
Variant & KB & n20 & n10 & n5 \\
\midrule
float32 (full) & 175.0 & 0.994\,/\,1.000 & 0.941\,/\,0.994 & 0.859\,/\,0.931 \\
uint8 (band) & 21.9  & 0.994\,/\,1.000 & 0.939\,/\,0.993 & 0.858\,/\,0.925 \\
uint4 (full) & 21.9  & 0.982\,/\,1.000 & 0.905\,/\,0.975 & 0.818\,/\,0.842 \\
uint4 (band) & 11.0  & 0.983\,/\,1.000 & 0.913\,/\,0.984 & 0.825\,/\,0.856 \\
\bottomrule
\end{tabular}
\end{minipage}\hfill
\begin{minipage}[t]{0.155\linewidth}
\centering
{\footnotesize (b) candidates\par}\vspace{2pt}
\setlength{\tabcolsep}{3pt}
\begin{tabular}{@{}c c c c@{}}
\toprule
$C$ & PESQ & n10 & n5 \\
\midrule
1  & 3.55 & 0.767 & 0.668 \\
4  & 3.53 & 0.889 & 0.785 \\
8  & 3.52 & 0.923 & 0.832 \\
16 & 3.51 & \textbf{0.944} & \textbf{0.861} \\
\bottomrule
\end{tabular}
\end{minipage}\hfill
\begin{minipage}[t]{0.35\linewidth}
\centering
{\footnotesize (c) desynchronization\par}\vspace{2pt}
\setlength{\tabcolsep}{3pt}
\begin{tabular}{@{}l c c@{}}
\toprule
Magnitude & pitch shift & speed \\
\midrule
small  & $\pm25$\,c: 1.000\,/\,1.000  & $3\%$: 0.959\,/\,0.997 \\
medium & $\pm50$\,c: 0.995\,/\,1.000  & $5\%$: 0.843\,/\,0.878 \\
large  & $\pm100$\,c: 0.784\,/\,0.728 & $10\%$: 0.675\,/\,0.290 \\
\bottomrule
\end{tabular}
\end{minipage}
\end{table*}


\textbf{Baselines.} We compare against (i) a re-implemented MelShield~\cite{jin2026}
and (ii) public AudioSeal~\cite{san2024} and WavMark~\cite{chen2023} models as
cross-paradigm waveform references (16 bits, post-vocoder). Our MelShield shares \texttt{RAWMER}'s frontend, band, payload and $\Delta_{\max}$,
uses mask floor $0.05$, energy exponent $0.75$, boundary margin $0.02$, zero headroom
and the calibrated $\tau = 0.61$ of~\cite{jin2026}; only $\alpha$ is tuned to match
clean PESQ ($0.050$ / $0.061$). MelShield reports ACC 1.000 under
non-noise attacks and 0.782/0.705 at 10/5~dB on HiFi-GAN, 0.779/0.701 on DiffWave, at
published points of PESQ\,$\approx$\,4.13 and 3.41~\cite{jin2026}. Ours reproduces
that signature and, on DiffWave at comparable clean PESQ (3.515 vs.\ 3.408), is
slightly \emph{stronger} than published (0.823 at 10~dB), so it is not weakened; on
HiFi-GAN our matched point gives 0.713 against a published 0.782 obtained in a
higher-quality regime with a differently scaled $\alpha$ ($0.25$, band $20\!:\!55$).
As $\alpha$ is not comparable across formulations, PESQ matching carries the whole
comparability burden; an ACC-vs-PESQ curve over an $\alpha$ sweep would remove it.

\subsection{Main Comparison}

Table~\ref{tab:main}a reports robustness at matched clean quality. On HiFi-GAN,
\texttt{RAWMER} lifts $\text{ACC}_{\text{noise10}}$ from 71.3\% to 94.1\% and
$\text{ACC}_{\text{noise5}}$ from 63.3\% to 85.9\%; on DiffWave, from 82.3\% to
92.2\% and 72.2\% to 83.5\%. Non-noise attacks stay at ACC\,$\geq\,$0.983 for all methods and are omitted for
space. Estimates come from one sampled set; a three-seed check (2026--2028, random 500)
varies by $\leq 0.002$ at 10 and 5\,dB, far below those gaps. Lacking
utterance-level confidence intervals, we treat small within-table gaps as
insignificant.

The waveform rows are context, not baselines to beat: they decode \emph{blindly}
while our detector reads the exact clean reference, and informed detection is strictly
easier~\cite{cox2008}. AudioSeal keeps moderate bit accuracy but rejects most noisy
samples (VR 0.14--0.30 at 20\,dB, zero at 10/5\,dB) and WavMark declines to decode
$\geq\!95\%$ of them---under heavy noise possibly correct blind behavior, not a
defect.

\textbf{Matched threshold.} \texttt{RAWMER}'s $\tau = 0.75$ is stricter than the
$\tau = 0.61$ calibrated in~\cite{jin2026}, so one may ask whether that margin alone
explains Table~\ref{tab:main}a; Table~\ref{tab:main}b reports VR under both.
MelShield's VR depends heavily on $\tau$: on HiFi-GAN at $\tau = 0.75$ it collapses to
0.434 at 10\,dB and 0.127 at 5\,dB, against \texttt{RAWMER}'s 0.991 and 0.918. The
relaxed threshold is not free---at chance accuracy the binomial false-positive rate is
0.35\% at $\tau = 0.75$ but 10.8\% at $\tau = 0.61$---so that column is a bit-margin
diagnostic, not a deployable operating point. We advocate $\tau = 0.75$ and use it for
all controls below.

\subsection{Reference-Assisted Verification}

\textbf{Negative controls.} On 2000 unattacked HiFi-GAN utterances with the correct
reference, key and payload, \texttt{RAWMER} passes at $\text{VR} = 1.000$; a wrong
key, payload or reference gives $\text{ACC} \in [0.499,0.503]$ and
$\text{VR} \in [0.0015,0.0040]$, matching the 0.35\% binomial tail---the observed
false positives are what chance alone predicts. Wrong-payload trials retain detector
\emph{confidence} on par with positives while their decoded bits disagree with the
claim: verification must hinge on payload match, not confidence.

\textbf{Reference compression.} The stored reference is the band $X_B$, not the full
mel. Table~\ref{tab:compress}a shows 8-bit quantization of $X_B$ shrinks it to about
22\,KB while ACC drops by at most 0.0013---$1/8$ of the float32 full-mel row
(175.0\,KB), or $1/4$ of the tighter float32 band-only baseline (87.5\,KB). The same
22\,KB on a 4-bit full-mel reference is consistently worse: precision inside the band
dominates coverage outside it. Under band-only storage our MelShield
re-implementation degrades far further (0.639/0.577/0.550 at 20/10/5\,dB versus
0.902/0.713/0.631 with a full float32 reference, random 500), as expected of an
absolute-residual detector.

\textbf{Blind partial-clip verification.} Because each bit repeats across keyed
blocks, \texttt{RAWMER} survives cropping. Without revealing the crop offset, a 25\%
slice from the start, middle or end is slid across the reference (stride 4 frames)
and the highest-scoring window taken: median start error 0--1 mel frames, ACC
0.985--0.991, $\text{VR} \geq 0.998$ at all three positions; 50\% clips reach
$1.000$.

\subsection{Ablation and Robustness Boundary}

\textbf{Which mechanism matters.} Table~\ref{tab:compress}b sweeps $C$: going from
$C = 1$ to $16$ lifts $\text{ACC}_{\text{noise5}}$ from 66.8\% to 86.1\% at negligible
PESQ cost, isolating Eq.~\eqref{eq:select} as the dominant contributor. Read honestly,
differential encoding \emph{alone} ($C = 1$, random pairs) beats quality-matched
MelShield at 5\,dB, 66.8\% vs.\ 63.3\%, but narrowly; selection supplies the remaining
19.3 points. We have not run the converse ablation---absolute-residual encoding
\emph{with} reliability-aware placement---so the differential score's independent
value rests on $C = 1$ alone.

\textbf{Robustness boundary.} \texttt{RAWMER} withstands 64~kbps MP3, 48~kbps AAC,
2~kHz lowpass, 8-bit quantization, $\pm 0.5$ clipping, 100~Hz--7~kHz bandpass and
stronger reverb ($\text{ACC} \geq 0.987$ on both vocoders), plus
EnCodec~\cite{defossez2022} at 6/12/24~kbps ($\text{ACC} \geq 0.976$ versus
0.888/0.942/0.954 for quality-matched MelShield on HiFi-GAN): the advantage is not
confined to Gaussian noise, though noise is where the methods separate most. At 0\,dB
SNR it degrades gracefully (ACC 0.73--0.75, $\text{VR} \leq 0.60$).

\textbf{Desynchronization is the attack to beat.} The genuine limit is
time--frequency desynchronization (Table~\ref{tab:compress}c). It arises from the
bounded-shift alignment of \S\ref{sec:method:det} breaking down, not from watermark
energy being lost---which is why it is dangerous. A keyless attacker who knows only that
alignment is frame-level can apply a $\pm 100$-cent shift or $\pm 10\%$ speed change,
both trivially scriptable and leaving speech intelligible, and will prefer this to
additive noise. Two keyless attacks are quantified in the supplement: [band-concentrated noise results TBD] and [re-vocoding results TBD]. Restoring
synchronization (multi-scale search, DTW) is thus a security requirement.
